\documentclass{article}
\usepackage{booktabs}

\title{Intermediate Value Performance Comparison}
\date{}

\begin{document}
\maketitle

\begin{table}[h]
\centering
\begin{tabular}{l r r r r r r r r r}
\toprule
 & \multicolumn{3}{c}{With intermediates}
 & \multicolumn{3}{c}{Without intermediates}
 & \multicolumn{3}{c}{Ratio (with / without)} \\
\cmidrule(lr){2-4} \cmidrule(lr){5-7} \cmidrule(lr){8-10}
Example
 & {Eval (ms)} & {Bwd (ms)} & {Vertices}
 & {Eval (ms)} & {Bwd (ms)} & {Vertices}
 & {Eval}      & {Bwd}      & {Vertices} \\
\midrule
conv.fld & 223.8 & 12.5 & 8737 & 206.9 & 9.0 & 8677 & 1.082 & 1.393 & 1.007 \\
\bottomrule
\end{tabular}
\caption{Comparison of backward-slicing performance with and without \texttt{@doc}-annotated intermediate values.
         Each row is an average over multiple runs. Ratio $> 1$ means the intermediate version is slower/larger.}
\label{tab:intermediate-perf}
\end{table}

\end{document}
